\section{Security}
\label{sec:security}

Papilot runs on the author's own computer and can be reached from other
devices over a private network. Its threat model is a web page in another tab
that tries to act on the author's behalf, and a device on the network that is
not the author's.

\subsection{Requests}

Every state-changing request must carry a token $\chi$ of $256$ random bits,
issued with the page and compared in constant time, which defends against
cross-site request forgery~\cite{barth2008}. A remote request must also come
from the configured origin and, over the private network, from the configured
owner. Writing $\rho$ for a request,
\begin{equation}
  \operatorname{allow}(\rho) =
  \bigl(\operatorname{safe}(\rho) \lor \operatorname{token}(\rho) = \chi\bigr)
  \land \operatorname{origin}(\rho) \in \mathcal{O}
  \land \operatorname{owner}(\rho) = o ,
  \label{eq:allow}
\end{equation}
where $\operatorname{safe}(\rho)$ holds for requests that only read. A guess
succeeds with probability $2^{-256}$, and a constant-time comparison leaks no
prefix of $\chi$ through timing.

\subsection{Agents}

Agents run with the author's own logins but within limits Papilot sets:
\begin{itemize}
  \item each task's agent works in its own worktree $w_i$ and is told never to
    edit the manuscript checkout $w_0$, other worktrees, or Papilot itself;
  \item commands run in the agent's sandbox, and anything that leaves it is
    either brought to the author for approval or, for an agent that cannot
    ask, declined in advance;
  \item only the author accepts changes into $w_0$, through the review of
    \cref{sec:model}.
\end{itemize}
The last point is the one that matters most: even an agent that ignored its
instructions could not change the manuscript without a decision the author
made while looking at the typeset result.
